\section{El estadístico y la regla de cierre}
\label{sec:estadistico}

\subsection{Para qué sirve un intervalo de confianza en esta campaña}

La objeción natural al intervalo de confianza es que añade ruido a la lectura. Es
la objeción correcta y la respuesta es que el intervalo no añade el ruido: lo
\emph{mide}. Sin él, el ruido sigue ahí y no se ve. Pero conviene ser preciso sobre
cuál es su función aquí, porque no es la función habitual en un artículo.

\textbf{El intervalo es un instrumento de selección, no de publicación.} Su
trabajo es decidir cuál de varias configuraciones avanza a la siguiente etapa. No
hace falta publicarlo para competir en una evaluación externa, donde el árbitro es
la tabla de resultados del organizador. Su valor está enteramente del lado
interno, y se concreta en cuatro cosas.

\subsubsection*{Primero: la cuantización hace que una estimación puntual decida
por debajo de lo que el dato resuelve}

Todo lo almacenado está redondeado a $10^{-4}$. El MDE de la celda peor resuelta es
0.0234. La razón es \textbf{234}. Una regla que compare dos máximos y se quede con
el mayor decide, en el caso peor, a partir de diferencias doscientas treinta veces
más pequeñas que lo que la población de esa celda puede distinguir de cero. Eso no
es una comparación estricta: es una moneda al aire con cuatro decimales.

\subsubsection*{Segundo: el número de oportunidades de acertar por casualidad es
grande, y crece con la ambición del diseño}

La canalización (D\ref{def:routing}) es la decisión donde esto muerde. Encaminar
configuraciones distintas a regiones distintas significa que el arm canalizado
\emph{elige una vez por región} mientras el arm base elige una sola vez. Con quince
regiones, el canalizado tiene quince oportunidades de parecer mejor por azar. Con
nueve paneles y varios arms, el número de contrastes llega al millar sin esfuerzo.

Esto no es una objeción a la canalización: es su condición de validez. El arm
canalizado gana con afinidad real \emph{cero} si no se le exige nada más. Lo único
que lo repara es el \textbf{ajuste cruzado}: la política de encaminamiento se ajusta
sobre unos pliegues y se evalúa sobre otros, de modo que la elección por región no
se puede premiar a sí misma. El intervalo es lo que permite comprobar que el
veredicto fuera de pliegue coincide con el veredicto dentro de pliegue; sin él, las
dos cifras son dos números que difieren y no hay forma de decir si la diferencia
importa.

\subsubsection*{Tercero: convierte \emph{esta celda no puede decir nada} en un
resultado escribible}

El vocabulario que la operación de comparación ya emite tiene seis palabras, y dos
de ellas son las que faltan en cualquier regla de estimación puntual:
\cod{null\_with\_power}, que es una celda con población suficiente cuyo intervalo
cubre el cero, y \cod{null\_unread}, que es una celda que no se pudo leer. La
distinción es la diferencia entre \emph{buscamos y no hay efecto} y \emph{no
pudimos buscar}, y una campaña que no la puede escribir tiene que convertir las dos
en un voto, en cualquier dirección, que es peor que abstenerse.

\subsubsection*{Cuarto: los envíos externos son un recurso escaso}

La validación final del proyecto es un despliegue en LAFA, y a futuro una
participación en CAFA. El número de configuraciones que se pueden enviar es
limitado y el ciclo de realimentación es de meses. Cada configuración enviada que
fue elegida sobre ruido es un envío desperdiciado y una espera perdida. La
validación externa no reduce la necesidad de rigor interno: la aumenta, porque
encarece el error.

\subsection{Qué estadístico decide}

Los dos candidatos no son dos formas de medir lo mismo:

\begin{center}
\begin{tabularx}{\linewidth}{@{}L{3.2cm} X X@{}}
\toprule
 & \cod{fmax\_w} & bootstrap pareado\\
\midrule
qué optimiza & una curva agregada sobre toda la población & el óptimo propio de cada proteína\\
qué premia   & predecir poco y con confianza & cubrir a cada proteína\\
admite $\sigma/\sqrt{n}$ & \textbf{no}. Es un máximo sobre $\tau$ de una combinación no lineal de dos medias & sí, por remuestreo de proteínas\\
comparable con CAFA & sí, es la medida del reto & no directamente\\
\bottomrule
\end{tabularx}
\end{center}

La consecuencia medida es que en los dos ejes que la campaña ya cerró los dos
estadísticos eligieron ganadores distintos, y en uno de los dos eligieron extremos
opuestos de la misma rejilla monótona. No hay respuesta independiente del
estadístico, y por tanto \textbf{el estadístico tiene que declararse, y entrar en
el sello}.

La recomendación es usar los dos con papeles separados y distintos:

\begin{itemize}[leftmargin=1.4em,itemsep=0.3em]
  \item \textbf{El bootstrap pareado decide} qué configuración avanza. Es la
        elección interna, y es la que necesita intervalo.
  \item \textbf{\cod{fmax\_w} se reporta} porque es la medida del reto y es lo que
        hace la cifra comparable con la literatura. Se publica como estimación
        puntual, sin intervalo inventado.
\end{itemize}

\subsection{La regla de cierre}

\begin{aviso}{La regla propuesta, y la que hay que sustituir}
La regla en vigor en el código compara el máximo de \cod{f\_micro\_w} del rival
contra el del incumbente, panel por panel, y exige unanimidad sobre los paneles que
testifican. Son estimaciones puntuales de una métrica cuantizada a $10^{-4}$, sin
intervalo. Además conviven tres reglas incompatibles: esa, una mayoría de los nueve
escrita en un plan, y el bootstrap con IC del 95\,\% escrito en las declaraciones.
Hay que quedarse con una.
\end{aviso}

La regla recomendada tiene tres cláusulas:

\begin{enumerate}[leftmargin=1.5em,itemsep=0.35em]
  \item \textbf{Una celda vota solo si su intervalo del 95\,\% contra el incumbente
        excluye el cero.} Una celda cuyo intervalo cubre el cero reporta
        \cod{null\_with\_power} o \cod{null\_unread} y \emph{no vota en ninguna
        dirección}.
  \item \textbf{Un eje sella solo si un nivel gana en todas las celdas que votan.}
  \item \textbf{Si no, el eje transporta el conjunto supervivente}
        (D\ref{def:haz}): los niveles que ninguna celda con voto llegó a batir.
\end{enumerate}

\subsection{El caso adverso, y por qué la regla necesita tres cláusulas más}

La regla anterior, tal como se enunció, tiene un agujero lógico y un modo de fallo
que conviene poner por escrito antes de adoptarla.

\textbf{El agujero.} Si ninguna celda vota, \emph{un nivel gana todas las celdas
que votan} es vacuamente cierto para todos los niveles a la vez. La regla no define
cuórum, de modo que en el régimen de efectos pequeños no sella: se rompe.

\textbf{El caso adverso, cuantificado con la tabla del propio documento.} Un efecto
uniforme de $+0{,}004$ solo es resoluble en \cod{PK:BPO}, cuyo MDE es $0{,}0032$;
por debajo de $0{,}0126$ ninguna celda \cod{NK} ni \cod{LK} puede votar. Es decir,
en ese régimen la regla degenera en \emph{decide \cod{PK:BPO}}, que es la categoría
\textbf{más fácil} y no la que el proyecto quiere titular. Peor: un efecto real y
pequeño con el mismo signo en 9 de 9 paneles y ningún intervalo excluyendo el cero
da cero votos y no sella, cuando una prueba de signos sobre nueve paneles da
$p = 0{,}004$. \textbf{La regla pierde una decisión correcta que la consistencia sí
sostiene}, y es exactamente la forma de la evidencia que la campaña ya produjo
(9 de 9 celdas, 36 de 36 contrastes).

\textbf{Y la multiplicidad corta en los dos sentidos.} El argumento a favor del
intervalo la invoca, y luego no se la aplica a sí mismo: al 95\,\% por celda, la
probabilidad de que al menos una de nueve vote en falso es del 37\,\%, y con
diecisiete regiones del 58\,\%. Bajo unanimidad, \emph{una sola celda falsa sella o
veta ella sola}. El ajuste cruzado arregla la selección, no la tasa de error por
familia.

Las tres cláusulas que faltan:

\begin{enumerate}[leftmargin=1.5em,itemsep=0.3em]
  \item \textbf{Cuórum.} Un eje no sella si votan menos de tres celdas. Con menos,
        el resultado es \emph{no resuelto}, que es distinto de \emph{no separó}.
  \item \textbf{Segunda vía por consistencia de signo.} Si ninguna celda vota pero
        las nueve tienen el mismo signo, el eje sella con una marca más débil, y la
        marca se publica: sellado por consistencia, no por separación.
  \item \textbf{Nivel por celda corregido.} Con nueve celdas, el nivel por celda
        sube al 99{,}4\,\% para mantener el 5\,\% por familia. Es conservador y hay
        que decir que lo es.
\end{enumerate}

Y el coste que el argumento a favor del intervalo no menciona: \textbf{no sellar
multiplica la rejilla del eje siguiente por el tamaño del haz}, hoy tres, y sin
sello no hay campeón que congelar en la etapa final. El argumento de que los envíos
externos son escasos vale igual en la dirección contraria.

\subsection{Resumen de la regla}

Arregla los dos modos de fallo a la vez: una celda sin poder no puede sostener una
mayoría ni puede vetarla. Es el estadístico que las declaraciones ya nombran y que
los 4{,}279 trabajos de comparación ya ejecutaron. Y es la única opción bajo la cual
\emph{el eje no eligió} es un resultado escribible, que es precisamente lo que uno
de los ejes de esta campaña produjo.

El coste hay que decirlo sin suavizarlo: esto vuelve las filas emparejadas por
proteína imprescindibles para cualquier sellado, y hoy viven solo en un registro de
eventos sin tipar y con borrado en cascada. Promoverlas a una tabla propia deja de
ser limpieza y se convierte en precondición de cerrar cualquier nodo.
